% !PRINTED 584
\weberpage{584}

To form the biquadratic equation whose roots are $\eta$,
$\eta_1$, $\eta_2$, $\eta_3$, we first seek the quadratic equation
with the roots $\eta$, $\eta_2$, and obtain it from

\[
\begin{aligned}
2(\eta+\eta_2)&=-1+\sqrt{n}\\
2(\eta-\eta_2)&=(i,\eta)+(-i,\eta).
\end{aligned}
\]

If we square these two equations and subtract the
second from the first, then by (30), (31) and (33) it follows that

\weq{35}{16\,\eta\,\eta_2=1+n-2\sqrt{n}-2n(-1)^f-2a(-1)^f\sqrt{n}.}

But by (29) and (35) the coefficients of the quadratic
equation sought are determined. It reads:
\[
\eta^2+\frac{\eta}{2}+\frac{1+n-2n(-1)^f}{16}=\left(\frac{\eta}{2}+\frac{a(-1)^f+1}{8}\right)\sqrt{n},
\]
and the biquadratic equation for the four periods is obtained
by squaring both sides:

\weq{36}{\left(\eta^2+\frac{\eta}{2}+\frac{1+n-2n(-1)^f}{16}\right)^2-n\left(\frac{\eta}{2}+\frac{a(-1)^f+1}{8}\right)^2=0.}

If we seek the coefficient of the first power of $\eta$, which
must be an integer, we find for it

\weq{37}{\frac{1+n-2n(-1)^f}{16}-n\,\frac{a(-1)^f+1}{8};}

$f$ is even when $n\equiv 1$, odd when $n\equiv 5$ (mod 8),
and accordingly
\[
1+n-2n(-1)^f\equiv 0\ (\mathrm{mod}\ 8).
\]
But from this it follows by (37) that $a(-1)^f+1$ must be
divisible by 4, hence

\weq{38}{a\equiv -(-1)^f\ (\mathrm{mod}\ 4),}

and from this it is further to be concluded, by (34), that $b$ must
be even, hence

\weq{39}{b\equiv 0\ (\mathrm{mod}\ 2).}

The biquadratic equation takes a simpler form under the substitution

\weq{40}{4\eta+1=\xi}

and by (36) becomes:

\weq{41}{[\xi^2+n(1-2(-1)^f)]^2-4n[\xi+(-1)^f a]^2=0,}

or, if we separate the two cases $n\equiv 1$ or $\equiv 5$ (mod 8):

\weq{42}{\begin{aligned}(\xi^2-n)^2-4n(\xi+a)^2&=0\\(\xi^2+3n)^2-4n(\xi-a)^2&=0.\end{aligned}}
